\documentclass{article}

\usepackage{amsmath} % math stuff
\usepackage{amssymb} % math stuff
\usepackage{array} % equations and stuff
\usepackage{bm} % bold math
%\usepackage{booktabs} % extra table rule options
%\usepackage{caption} % suppressed table numbering; incompatible with revtex, and longtable, I think
\usepackage{comment} % comment environment
%\usepackage{enumitem} % customization of enumeration, itemize, and description
\usepackage[T1]{fontenc} % font encoding for special characters, must also use scalable font package
\usepackage[margin=0.8in]{geometry} % paper sizes and margins (but be careful not to mess up pre-defined pages)
\usepackage{graphicx} % for graphics
%\usepackage{helvet} % default font is the helvetica postscript font
\usepackage[utf8]{inputenc} % special characters in tex input
\usepackage{layouts} % print units like widths
\usepackage{lipsum} % lorem ipsum filler text
\usepackage{lmodern} % scalable font?
\usepackage{longtable} % multi-page tables
\usepackage{makecell} % specify line-breaks in table cells
\usepackage{mathrsfs} % math script font
\usepackage{mhchem} % easier chemical formula
\usepackage{microtype} % allows disabling of ligatures
\usepackage{multicol} % multicolumns
%\usepackage{newcent} % new century schoolbook font
\usepackage{nicefrac}
\usepackage{numprint} % print and format (large) numbers
\usepackage{parskip} % removes paragraph indentation, and adjusts paragraph skip, as well as list items
\usepackage{pdfpages} % add pdf files as pages
%\usepackage{setspace} % adjust text spacing and indents
\usepackage{siunitx} % decimal alignment, and degree symbol
\usepackage{subfigure} % divided figures
%\usepackage{tabu} % extra table options
\usepackage{textcomp} % symbols
\usepackage{threeparttablex} % better footnotes with longtable
\usepackage{titling} % title placement
\usepackage{ulem} % strikethrough text
\usepackage{upgreek} % upright Greek letters
%\usepackage{url} % superceded by hyperref
\usepackage{verbatim} % verbatim environment
\usepackage{xcolor} % colors and color boxes
\usepackage{xspace} % commands that don't eat up white space
\usepackage{hyperref} % links and page setup; should always come last

\hypersetup{
 bookmarks=true,
 colorlinks=true,
 citecolor=blue,
 linkcolor=blue,
 urlcolor=blue,
 pdfstartview={XYZ null null 1.0} % default open view is 100%
}

\DisableLigatures[f,t]{encoding = T1} % disable ff, fi, fl, tt ligatures; without options, it also disables -- = endash
\renewcommand{\arraystretch}{1.0} % extra vertical (and horizontal?) space in tables

% define centered, left- and right-aligned columns with specified widths
\newcommand{\PreserveBackslash}[1]{\let\temp=\\#1\let\\=\temp}
\newcolumntype{C}[1]{>{\PreserveBackslash\centering}p{#1}}
\newcolumntype{L}[1]{>{\PreserveBackslash\raggedright}p{#1}}
\newcolumntype{R}[1]{>{\PreserveBackslash\raggedleft}p{#1}}

\begin{document}

\pagestyle{empty} % don't number pages

% custom title
\begin{center}
{\LARGE Classic Riddler}

\vspace{0.15in}

{\Large 28 January 2022}
\end{center}


\section*{Riddle:}

You want to change the transmission fluid in your old van, which holds 12 quarts of fluid.
At the moment, all 12 quarts are ``old.''
But changing all 12 quarts at once carries a risk of transmission failure.

Instead, you decide to replace the fluid a little bit at a time.
Each month, you remove one quart of old fluid, add one quart of fresh fluid and then drive the van to thoroughly mix up the fluid.
(I have no idea if this is mechanically sound, but I'll take Travis's word on this!)
Unfortunately, after precisely one year of use, what was once fresh transmission fluid officially turns ``old.''

You keep up this process for many, many years.
One day, immediately after replacing a quart of fluid, you decide to check your transmission.
What percent of the fluid is old?


\section*{Solution:}

At the moment you add in new fluid, \nicefrac{1}{12} of the fluid is brand new.
Since one month prior, you made the same change, where \nicefrac{1}{12} of that month's fluid was brand new, of the remaining \nicefrac{11}{12}, \nicefrac{1}{12} is one month old, or $\nicefrac{11}{12}\times\nicefrac{1}{12}$ total.
Continuing to the next previous month, $(\nicefrac{11}{12})^2\times\nicefrac{1}{12}$ of the fluid is two months old.
For the last twelve additions of fluid, the total that is not yet old can be written as a sum:

\[
\sum_{n=0}^{11}\left(\frac{1}{12}\right)\left(\frac{11}{12}\right)^n
\]

This sums to approximately 0.648.
The remainder of the fluid is at least twelve months old, so the solution is (approximately) \fcolorbox{red}{white}{\textbf{0.352}}\,.
As the frequency of fluid changes increases (to say, daily or more), this value approaches $\nicefrac{1}{e}\approx0.368$.


\end{document}